\section{Разработка программного решения для анализа массива данных}

\AKHworries{АХ: для описания алгоритмов используем  algorithm2e с параметрами linesnumbered,ruled,vlined. При этом внутри псевдокода все на английском. Описание алгоритма технически корректным языком. Третируем русский и английский текст.  }
\subsection{Алгоритм обратного восстановления исходного набора коробок}

В данном разделе представлен алгоритм обратного восстановления исходного набора коробок по результату их размещения на паллете. В качестве входных данных алгоритм принимает файл, содержащий идентификатор коробки (sku) и координаты ее расположения на паллете (левый нижний угол и правый верхний). Также предусмотрено добавление дополнительных атрибутов: вес и тип коробки. Псевдокод алгоритма приведён в алгоритме~\ref{alg:reverse_reconstruct}.

Геометрические параметры коробок восстанавливаются по разностям координат вдоль осей $x$, $y$ и $z$, что при предположении ортогональной ориентации и отсутствии деформаций однозначно определяет длину, ширину и высоту. Все вычисления проводятся в целочисленном формате, ввиду предположения, что все координаты являются целыми числами.

Алгоритм реализует два режима работы. В базовом режиме каждая размещённая коробка обрабатывается как отдельный объект, что обеспечивает прямую обратимость преобразования. В режиме агрегации выполняется объединение идентичных коробок с формированием компактного представления исходного заказа.

Агрегация коробок производится только в случае полного совпадения всех атрибутов. Для эффективного выявления идентичных коробок используется хэш-таблица с составным ключом, процедура вычисления которого описана в алгоритме~\ref{alg:reverse_key_hash}. Поскольку вес коробок задан в граммах и является целочисленным, он включается в ключ без дополнительного преобразования.

Для каждой коробки вычисляется соответствующий ключ и выполняется проверка его наличия в хэш-таблице. При отсутствии ключа создаётся новая запись, в противном случае увеличивается счётчик коробок в существующей записи. Использование хэш-таблицы обеспечивает амортизированное константное время операций поиска и вставки, что приводит к линейной временной сложности алгоритма по числу входных данных.


\begin{algorithm}[H]
\caption{Reconstruct boxes from a pallet placement}
\label{alg:reverse_reconstruct}
\KwIn{$placement$ -- list of placed records; $useMerge$ -- aggregation flag}
\KwOut{$boxes$ -- reconstructed list of box records}

$boxes \leftarrow \emptyset$\;
\If{$useMerge$}{
    $index \leftarrow$ empty hash table\;
}

\ForEach{$r \in placement$}{
    $(L,W,H) \leftarrow$ \textnormal{ExtractDimensions}$(r)$\;
    $(weight,aisle) \leftarrow$ \textnormal{ExtractAttributes}$(r)$\;

    \If{not $useMerge$}{
        append $(r.sku, L, W, H, weight, aisle, 1)$ to $boxes$\;
    }
    \Else{
        $key \leftarrow (r.sku, L, W, H, aisle, weight)$\;
        \If{$key \notin index$}{
            append $(r.sku, L, W, H, weight, aisle, 1)$ to $boxes$\;
            $index[key] \leftarrow |boxes|-1$\;
        }
        \Else{
            $j \leftarrow index[key]$\;
            $boxes[j].quantity \leftarrow boxes[j].quantity + 1$\;
        }
    }
}

\Return{$boxes$}\;
\end{algorithm}



\begin{algorithm}[H]
\caption{Hash computation for a composite key}
\label{alg:reverse_key_hash}
\KwIn{$key$ -- tuple $(sku, L, W, H, aisle, weight)$}
\KwOut{$h$ -- hash value of type \texttt{size\_t}}
$h \leftarrow 1469598103934665603$\;
\ForEach{$v \in [sku, L, W, H, aisle, weight]$}{
    $h \leftarrow h \oplus \Big( cast(v) + 0x9b97f4a7c15 + (h \ll 6) + (h \gg 2) \Big)$\;
}
\Return{$h$}\;
\end{algorithm}

\subsection{Программный анализ геометрии и устойчивости паллетной укладки}

В рамках проекта разработан программный модуль анализа результатов паллетной укладки, основанный на геометрическом представлении размещённых коробок. Общая схема работы модуля приведена в алгоритме~\ref{alg:pallet_analysis}. На вход модуль принимает набор коробок с заданными координатами размещения и геометрические параметры основания паллеты, формируя совокупность количественных метрик, характеризующих качество и устойчивость укладки.

На первом этапе выполняется вычисление глобальных характеристик паллеты: итоговой высоты, суммарного объёма груза и коэффициента объёмного заполнения. Одновременно проводятся проверки геометрической корректности размещения, включая выход коробок за границы основания паллеты и наличие некорректных размеров. Кроме того, вычисляется центр масс паллеты в горизонтальной плоскости на основе координат центров коробок и их весов.

На следующем этапе анализируется корректность взаимного расположения коробок. Для каждой пары коробок проверяется наличие пересечения в трёхмерном пространстве, что позволяет выявлять недопустимые конфигурации укладки. Данная процедура имеет квадратичную вычислительную сложность $O(n^2)$, где $n$ — число коробок.

Ключевым этапом анализа является вычисление коэффициента опоры коробок (алгоритм~\ref{alg:support_ratio}). Для коробок, размещённых непосредственно на основании паллеты, коэффициент опоры полагается равным единице. Для остальных коробок определяется множество опорных коробок нижнего уровня, находящихся с ними в контакте по высоте. Площадь опоры вычисляется как площадь объединения областей пересечения верхних проекций опорных коробок с основанием анализируемой коробки.

Задача вычисления площади опоры сводится к нахождению площади объединения прямоугольников на плоскости. Для этого область контакта разбивается по одной из горизонтальных координат, после чего для каждого интервала вычисляется длина объединения отрезков по второй координате, а итоговая площадь определяется суммированием вкладов всех интервалов.

Вычисление коэффициентов опоры для всех коробок имеет сложность $O(n \cdot k^2)$, где $n$ — общее число коробок, а $k \ll n$ — среднее число опорных коробок. На основе полученных значений формируются агрегированные метрики устойчивости укладки, включая среднее и минимальное значения коэффициента опоры, долю коробок, удовлетворяющих заданному порогу устойчивости, а также число коробок, не имеющих опоры при ненулевой высоте размещения.

Разработанный модуль анализа позволяет перейти от геометрического описания паллетной укладки к формализованному набору количественных показателей, что обеспечивает автоматизированное сравнение алгоритмов паллетизации и использование модуля в составе вычислительных экспериментов.

\begin{algorithm}[H]
\caption{Support ratio computation for a box}
\label{alg:support_ratio}
\KwIn{$A$ -- target box; $P$ -- placed boxes}
\KwOut{$r$ -- support ratio}

\If{$A.z_1 = 0$}{
    \Return{$1$}
}

$R \leftarrow$ find boxes $B \in P$ such that $B.z_2 = A.z_1$\;
$C \leftarrow$ compute contact areas between $A$ and boxes in $R$\;
$supportedArea \leftarrow$ \textnormal{UnionContactArea}$(C)$\;

$r \leftarrow supportedArea/\textnormal{BaseArea}(A)$\;
$r \leftarrow \min(1,\max(0,r))$\;

\Return{$r$}\;
\end{algorithm}

\begin{algorithm}[H]
\caption{Pallet analysis and metric computation}
\label{alg:pallet_analysis}
\KwIn{$P$ -- list of placed boxes; $L,W$ -- pallet base dimensions; $s_{\min}$ -- support threshold}
\KwOut{$M$ -- pallet-level metrics; $support[i]$ -- per-box support ratios}

$maxHeight \leftarrow 0$\;
$totalVolume \leftarrow 0$\;

\ForEach{$b \in P$}{
    $maxHeight \leftarrow \max(maxHeight, b.z_2)$\;
    $totalVolume \leftarrow totalVolume + \textnormal{Volume}(b)$\;
}

$overlapPairs \leftarrow$ \textnormal{CountOverlaps3D}$(P)$\;

$floating \leftarrow 0$; $goodSupport \leftarrow 0$\;
\For{$i \leftarrow 1$ \KwTo $|P|$}{
    $support[i] \leftarrow$ \textnormal{SupportRatio}$(P[i], P)$\;
    \If{$P[i].z_1 > 0$ \textbf{and} $support[i] = 0$}{
        $floating \leftarrow floating + 1$\;
    }
    \If{$support[i] \ge s_{\min}$}{
        $goodSupport \leftarrow goodSupport + 1$\;
    }
}

$M.maxHeight \leftarrow maxHeight$\;
$M.fillRatio \leftarrow totalVolume/(L\cdot W\cdot maxHeight)$\;
$M.goodSupportFraction \leftarrow goodSupport/|P|$\;
$M.floatingBoxes \leftarrow floating$\;
$M.overlapPairs \leftarrow overlapPairs$\;

\Return{$M, support$}\;
\end{algorithm}



\begin{algorithm}[H]
\caption{Union area computation of contact regions}
\label{alg:union_area}
\KwIn{$C$ -- set of rectangular contact regions}
\KwOut{$A$ -- total contact area}

$A \leftarrow 0$\;
split $C$ into disjoint vertical intervals\;
\ForEach{interval}{
    compute union length along the second coordinate\;
    add interval contribution to $A$\;
}
\Return{$A$}\;
\end{algorithm}

\subsection{Построение алгоритма укладки паллеты}
В качестве базового решения в работе рассматривается слоевой жадный алгоритм укладки паллеты, основанный на эвристике \textit{Candidate Heights}, формально описанной в алгоритме~\ref{alg:candidate_heights}. Алгоритм используется как тестовое решение для оценки производительности и качества укладки, а также для последующего сравнения с более сложными методами.

Алгоритм реализует слоевой подход: на каждом шаге выбирается фиксированная высота слоя, после чего выполняется двумерная укладка коробок в его пределах с использованием эвристики \textit{MaxRects} с критерием максимального заполнения площади и разрешением поворота коробок в плоскости основания. Коробки обрабатываются в фиксированном порядке, определяемом площадью основания и высотой.

Для определения высоты слоя формируется множество потенциальных значений на основе высот неразмещённых коробок. Для каждого кандидата выполняется пробная двумерная укладка коробок с высотой, не превышающей выбранное значение кандидата. Для каждой укладки вычисляется коэффициент объёмного заполнения. В качестве итоговой высоты слоя выбирается значение, соответствующее максимальному коэффициенту заполнения.

После выбора высоты производится фактическая укладка коробок с фиксацией координат, после чего алгоритм переходит к следующему слою и повторяет процедуру до размещения всех коробок либо невозможности дальнейшей укладки.

Рассматриваемый алгоритм не учитывает условия устойчивости укладки, включая площадь опоры, характер контакта с нижележащими коробками и положение центра масс. В связи с этим он ориентирован исключительно на геометрическое заполнение объёма и используется в работе в качестве базового сравнительного метода.

В таблице~\ref{tab:candidate_heights_results} приведены экспериментальные результаты работы алгоритма в зависимости от числа кандидатных высот $k$, включая среднее время выполнения и коэффициент объёмного заполнения паллеты.

\begin{algorithm}[H]
\caption{Layer-based palletization with Candidate Heights}
\label{alg:candidate_heights}
\KwIn{$B$ -- list of boxes; $K$ -- number of candidate heights}
\KwOut{$placement$ -- list of box placements}

sort $B$ by decreasing base area and height\;
initialize empty pallet $P$\;

\While{there exist unpacked boxes}{
    $C \leftarrow$ \textnormal{SelectTopKHeights}$(B, K)$\;
    $H^\star \leftarrow$ height from $C$ maximizing \textnormal{EvaluateLayerFill}$(B, H)$\;

    start new layer of height $H^\star$\;
    \ForEach{box $b \in B$}{
        \If{$b$ is unpacked \textbf{and} $b.height \le H^\star$}{
            \If{\textnormal{FitsInLayer}$(P, b)$}{
                place $b$ in current layer using MaxRects\;
                mark $b$ as packed\;
            }
        }
    }
}

\Return{$placement$}\;
\end{algorithm}


\begin{table}[H]
\centering
\caption{Performance of the Candidate Heights algorithm for different values of $k$}
\label{tab:candidate_heights_results}
\begin{tabular}{|c|c|c|}
\hline
Number of candidates $k$ & Average time (ms) & Percolation \\ \hline
10 & 67 & 0.7655 \\ \hline
50 & 277 & 0.7866 \\ \hline
100 & 536 & 0.7944 \\ \hline
200 & 1006 & 0.8009 \\ \hline
500 & 1134 & 0.8025 \\ \hline
\end{tabular}
\end{table}

На основе базового слоевого подхода был рассмотрен альтернативный алгоритм \emph{height-map palletization}, в котором поверхность паллеты представляется в виде разбиения основания на непересекающиеся прямоугольные плитки (tiles), каждая из которых характеризуется собственной высотой $z$ и флагом блокировки. В отличие от классической слоевой схемы, данный подход явно моделирует неоднородную поверхность и позволяет целенаправленно обрабатывать локальные области с пониженной высотой.

На каждой итерации определяется текущий уровень слоя $L_0$ как максимальная высота среди незаблокированных плиток, после чего выполняется двухэтапная обработка поверхности. Сначала заполняются плитки с $z < L_0$ («провалы»): для выбранной плитки подбирается коробка, полностью помещающаяся в её область (с возможным поворотом). В зависимости от доли покрытия основания плитки коробка либо размещается с подъёмом всей плитки, либо выполняется локальное размещение с ограниченным числом попыток, после чего оставшиеся области могут быть заблокированы. Такой механизм предотвращает зацикливание на сильно фрагментированных участках.

После обработки «провалов» выполняется укладка на верхних областях (плитки на высоте $z = L_0$). На этих регионах применяется двумерная упаковка в сочетании с эвристикой \textit{Candidate Heights}, что позволяет выбирать высоту, наиболее эффективную для текущей конфигурации поверхности. Если верхняя поверхность оказывается слишком раздробленной и размещение становится невозможным, выполняется операция \emph{flattening}, при которой все доступные плитки поднимаются до общего уровня и сливаются, после чего укладка повторяется.

Основным недостатком алгоритма является то, что он не гарантирует размещение всех коробок. По мере формирования неоднородной поверхности и последовательного разбиения плиток часть оставшихся коробок перестаёт помещаться внутрь доступных прямоугольных областей. В результате алгоритм завершается преждевременно с неполной упаковкой, что ограничивает его применимость для некоторых заказов. Псевдокод алгоритма представлен в алгоритме~\ref{alg:heightmap_palletization}.

\begin{algorithm}[H]
\caption{Height-map palletization}
\label{alg:heightmap_palletization}
\KwIn{$B$ -- list of boxes; pallet base $(L,W)$}
\KwOut{$P$ -- list of placed boxes}

sort $B$ by decreasing base area and height\;
$T \leftarrow \{(0,0,L,W,z{=}0,\textnormal{blocked}{=}\textnormal{false})\}$\;
$P \leftarrow \emptyset$\;

\While{there exist unpacked boxes}{
    $L_0 \leftarrow \max\{t.z \mid t\in T,\ \neg t.\textnormal{blocked}\}$\;

    \While{exists non-blocked tile $t\in T$ with $t.z<L_0$}{
        place one fitting box inside $t$ (rotation allowed)\;
        update $T$ by splitting and changing tile heights\;
        add placed box to $P$\;
        \If{no box fits in any hole tile}{
            mark remaining hole tiles as \textnormal{blocked}\;
            \textbf{break}\;
        }
    }

    $R \leftarrow \{t \in T \mid \neg t.\textnormal{blocked} \ \wedge\ t.z=L_0\}$\;
    $H^\star \leftarrow \arg\max_H \textnormal{EvalHeight}(B,R,H)$\;
    $placed \leftarrow \textnormal{MaxRectsPackAndUpdateTiles}(B,T,R,L_0,H^\star)$\;

    \If{$\neg placed$}{
        flatten all non-blocked tiles to the current maximum height and merge\;
        $L_0 \leftarrow \max\{t.z \mid t\in T,\ \neg t.\textnormal{blocked}\}$\;
        $R \leftarrow \{t \in T \mid \neg t.\textnormal{blocked} \ \wedge\ t.z=L_0\}$\;
        $placed \leftarrow \textnormal{MaxRectsPackAndUpdateTiles}(B,T,R,L_0,H^\star)$\;
    }

    \If{$\neg placed$}{\textbf{break}\;}
}
\Return{$P$}\;
\end{algorithm}


В качестве альтернативы слоевому подходу был разработан алгоритм укладки, основанный на размещении коробок в трёхмерном пространстве по набору кандидатных позиций. При каждой успешной постановке коробки определяется ограниченное множество перспективных точек размещения, соответствующих правым и верхним граням уже размещённых коробок по осям $x$, $y$ и $z$. Такой подход существенно сокращает пространство поиска по сравнению с непрерывным перебором и при этом сохраняет возможность формирования плотных укладок.

В начале алгоритма производится упорядочивание коробок по их площади основания и высоте. В процессе алгоритма коробки обрабатываются в данном порядке. Для каждой коробки перебираются две ориентации в плоскости основания (поворот на $90^\circ$) и набор кандидатных позиций. Каждая кандидатная постановка проверяется на корректность: коробка должна полностью помещаться в пределах основания паллеты, не пересекаться с уже размещёнными коробками и удовлетворять ограничению устойчивости.

Проверка устойчивости реализована эвристически. Для коробок, размещаемых на основании паллеты, дополнительных условий не требуется. Для коробок, устанавливаемых выше, оценивается доля площади опоры как отношение площади пересечения основания новой коробки с верхними гранями коробок нижнего уровня. Дополнительно проверяется, что проекция центра основания, используемая как приближение центра масс, лежит в области контакта хотя бы с одной опорной коробкой. В финальном критерии допускается малая опора (не менее $10\%$) при обязательной поддержке центра масс.

Выбор наилучшей позиции осуществляется с использованием лексикографической функции качества, минимизирующей высоту размещения $z$ и затем координаты $x$ и $y$, с дополнительным поощрением большей доли опоры. После размещения коробки множество кандидатных позиций расширяется точками, соответствующими её граням, а также, при ограниченном размере пространства, комбинациями ранее встречавшихся значений координат. Такой механизм обеспечивает компромисс между качеством поиска и вычислительной сложностью алгоритма.

На рисунке~\ref{rs_153.png} показан пример визуализации результата укладки, полученного описанным алгоритмом, иллюстрирующий структуру слоёв и распределение коробок по высоте паллеты.

\begin{algorithm}[H]
\caption{3D candidate-position palletization with stability check}
\label{alg:candidate_positions_3d}
\KwIn{$B$ -- list of box specifications; pallet base $(L,W)$}
\KwOut{$P$ -- list of placed boxes with coordinates}

expand quantities in $B$ into individual boxes; sort $B$ by decreasing base area, weight, height\;
$P \leftarrow \emptyset$, $C \leftarrow \{(0,0,0)\}$ \tcp*{placed boxes, candidate positions}

\ForEach{box $b \in B$}{
    $best \leftarrow \textnormal{None}$, $bestScore \leftarrow +\infty$\;
    $C' \leftarrow \textnormal{BuildCandidateSet}(C,P)$\;

    \ForEach{$(o_L,o_W) \in \{(b.L,b.W),(b.W,b.L)\}$}{
        \ForEach{$(x,y,z) \in C'$}{
            $(X,Y,Z) \leftarrow (x+o_L,\; y+o_W,\; z+b.H)$\;
            \If{not \textnormal{InsidePallet}$(x,y,X,Y)$
                or \textnormal{IntersectsAny}$(P,x,y,z,X,Y,Z)$
                or not \textnormal{StableOnSupport}$(P,x,y,z,X,Y)$}{continue}

            $s \leftarrow \textnormal{SupportRatio}(P,x,y,z,X,Y)$\;
            \If{\textnormal{PlacementScore}(x,y,z,s) < bestScore}{
                $best \leftarrow (b.sku,x,y,z,X,Y,Z,b.weight)$\;
            }
        }
    }

    \If{$best \neq \textnormal{None}$}{
        append $best$ to $P$; update $C \leftarrow C \cup \textnormal{NewCandidatesFromPlacement}(best)$\;
    }
}
\Return{$P$}\;
\end{algorithm}

\begin{figure}[H]
\centering
\includegraphics[width=0.7\textwidth]{figures/rs_153.png}
\caption{Пример визуализации результата укладки, полученного алгоритмом 3D candidate-position}
\label{rs_153.png}
\end{figure}


\AKHworries{АХ:  алгоритм Романа.  }